\subsection{GROUP DEFINED BY A COMPLETE NORMED LIE ALGEBRA}

Let \(H = \sum_{r,s \geq 0} H_{r,s} \in \hat{L}(\{U, V\})\) be the Hausdorff series (§ 6, no. 4, Definition 1). We shall show that the corresponding formal power series with continuous components

\[
\tilde {\mathrm{H}} = \sum_ {r, s \geqslant 0} \tilde {\mathrm{H}} _ {r, s} \in \hat {\mathrm{P}} (\mathfrak {g} \times \mathfrak {g}, \mathfrak {g})\tag{10}
\]

is convergent (Differentiable and Analytic Manifolds, R, 4.1.1).

Let \(r \geqslant 0\) , \(s \geqslant 0\) be such that \(r + s \neq 0\) and let \(\| \tilde{\mathrm{H}}_{r,s} \|\) be the norm of the continuous polynomial \(\tilde{\mathrm{H}}_{r,s}\) (Differentiable and Analytic Manifolds, R, Appendix, no. 2).

Lemma 3.

\[
\left\| \tilde {\mathrm{H}} _ {r, s} \right\| \leqslant a ^ {- (r + s - 1) \theta}.
\]

Let B be a Hall set relative to I and let \(H_{r,s} = \sum_{b \in B} \lambda_b e_b\) be the decomposition of \(H_{r,s}\) with respect to the corresponding basis of \(L(\{U, V\})\) . Then

\[
\left| \lambda_ {b} \right| \leqslant a ^ {- (r + s - 1) \theta}.\tag{11}
\]

This is trivial for \(p = 0\) , for \(\lambda_b \in \mathbf{Q}\) ; and it follows from Proposition 1 of no. 1 for \(p \neq 0\) .

Moreover,

\[
\left\| \tilde {e} _ {b} \right\| \leqslant 1 \quad \text { for } b \in \mathbf {B}.\tag{12}
\]

We show more generally by induction on \(n\) that, for every alternant \(b\) of degree \(n\) in the two indeterminates U and V (§2, no. 6), \(\| \tilde{b} \| \leqslant 1\) . If \(n = 1\) , \(\tilde{b}\) is one of the projections of \(\mathfrak{g} \times \mathfrak{g}\) onto \(\mathfrak{g}\) and hence is of norm \(\leqslant 1\) ; if \(n > 1\) , there exist two alternants \(b_1\) and \(b_2\) of degrees \(< n\) such that \(b = [b_1, b_2]\) . As the mapping \(\gamma: (x, y) \mapsto [x, y]\) of \(\mathfrak{g} \times \mathfrak{g}\) into \(\mathfrak{g}\) is bilinear of norm \(\leqslant 1\) , we have (Differentiable and Analytic Manifolds, R, Appendix, no. 4)

\[
\| \tilde {b} \| = \| \gamma \circ (\tilde {b} _ {1}, \tilde {b} _ {2}) \| \leqslant \| \tilde {b} _ {1} \|. \| \tilde {b} _ {2} \| \leqslant 1.\tag{13}
\]

Relations (11) and (12) imply the lemma.

PROPOSITION 2. The formal power series \(\tilde{\mathbf{H}}\) is a convergent series (Differentiable and Analytic Manifolds, R, 4.1.1). If \(\mathbf{G}\) is the ball \(\{x\in\mathfrak{g}\mid\|x\|<a^{\theta}\}\) , the domain of absolute convergence of \(\tilde{\mathbf{H}}\) (Differentiable and Analytic Manifolds, R, 4.1.3) contains \(\mathbf{G}\times\mathbf{G}\) .

If \(u\) and \(v\) are two real numbers \(>0\) such that \(u < a^{\theta}\) and \(v < a^{\theta}\) , then (Lemma 3)

\[
\| \tilde {\mathrm{H}} _ {r, s} \| u ^ {r} v ^ {s} \leqslant a ^ {\theta} (u a ^ {- \theta}) ^ {r} (v a ^ {- \theta}) ^ {s}\tag{14}
\]

and \(\| \tilde{\mathrm{H}}_{r,s}\| u^r v^s\) tends to 0 when \(r + s\) tends to infinity.

We denote by \(h \colon G \times G \to \mathfrak{g}\) the analytic function (Differentiable and Analytic Manifolds, R, 4.2.4) defined by \(\tilde{\mathbf{H}}\) , that is by the formula

\[
h (x, y) = \sum_ {r, s \geqslant 0} \tilde {\mathrm{H}} _ {r, s} (x, y) = \sum_ {r, s \geqslant 0} \mathrm{H} _ {r, s} (x, y) \quad \text { for } (x, y) \in \mathrm{G} \times \mathrm{G}.\tag{15}
\]

This function is called the Hausdorff function of g.

Let \((x,y)\in \mathbf{G}\times \mathbf{G}\) . Then

\begin{enumerate}
\def\labelenumi{(\arabic{enumi})}
\setcounter{enumi}{15}
\tightlist
\item
\end{enumerate}

\[
\left\| \tilde {\mathrm{H}} _ {r, s} (x, y) \right\| \leqslant \sup (\left\| x \right\|, \left\| y \right\|)\tag{16}
\]

\[
\| h (x, y) \| \leqslant \sup (\| x \|, \| y \|).\tag{17}
\]

\begin{enumerate}
\def\labelenumi{(\arabic{enumi})}
\setcounter{enumi}{16}
\tightlist
\item
  follows immediately from (16) and (16) is trivial for \(r = s = 0\) ; if \(r \geqslant 1\) , then
\end{enumerate}

\[
\begin{array}{r l} \| \tilde {\mathrm{H}} _ {r, s} (x, y) \| & \leqslant \| \tilde {\mathrm{H}} _ {r, s} \| \| x \| ^ {r} \| y \| ^ {s} \\ & \leqslant \| x \| \left(\frac {\| x \|}{a ^ {\theta}}\right) ^ {r - 1} \left(\frac {\| y \|}{a ^ {\theta}}\right) ^ {s} \\ & \leqslant \| x \|; \end{array}
\]

we argue similarly if \(s \geqslant 1\) .

In particular, \(\| h(x,y)\| < a^{\theta}\) for \((x,y)\in \mathbf{G}\times \mathbf{G}\) .

PROPOSITION 3. Let G be the ball \(\{x\in\mathfrak{g}\mid\|x\|<a^{\theta}\}\) . The analytic mapping

\[
h \colon \mathrm{G} \times \mathrm{G} \rightarrow \mathrm{G}
\]

makes G into a group in which 0 is the identity element and -x is the inverse of x for all \(x \in G\) . Moreover, if R is a real number such that \(0 < R < a^{\theta}\) , the ball

\[
\{x \in \mathfrak {g} | \| x \| <   \mathrm{R} \}
\]

(resp. \(\{x \in \mathfrak{g} \mid \|x\| \leqslant R\}\)) is an open subgroup of G.

As \(\mathrm{H}(\mathrm{U}, -\mathrm{U}) = 0\) and \(\mathrm{H}(0, \mathrm{U}) = \mathrm{H}(\mathrm{U}, 0) = \mathrm{U}\) , \(h(x, -x) = 0\) and

\[
h (0, x) = h (x, 0) = x
\]

for all \(x \in G\) . It therefore remains to prove the associativity formula

\[
h (h (x, y), z) = h (x, h (y, z)) \quad \text { for } x, y, z \text { in } G.\tag{18}
\]

As

\[
\mathrm{H} (\mathrm{H} (\mathrm{U}, \mathrm{V}), \mathrm{W}) = \mathrm{H} (\mathrm{U}, \mathrm{H} (\mathrm{V}, \mathrm{W}))
\]

in \(\hat{\mathbf{L}} (\{\mathrm{U},\mathrm{V},\mathrm{W}\})\) (§ 6, no. 5, Proposition 4), we have

\[
\tilde {\mathrm{H}} \circ (\tilde {\mathrm{H}} \times \mathrm{Id} _ {\mathfrak {g}}) = \tilde {\mathrm{H}} \circ (\mathrm{Id} _ {\mathfrak {g}} \times \tilde {\mathrm{H}})\tag{19}
\]

in \(\hat{\mathbf{P}} (\mathfrak{g}\times \mathfrak{g}\times \mathfrak{g};\mathfrak{g})\) (no. 2) and (19) implies (18) by (16) and Differentiable and Analytic Manifolds, R, 4.1.5.

*In other words (Chapter III, §1), G with the Hausdorff function is a Lie group.*

